% Three levels: both palettes, longer captions, explicit breaks and editable icons.
\ifdefined\pdfminorversion\pdfminorversion=7\fi
\documentclass[aspectratio=169]{beamer}
\usetheme{bitpolito-studio}
\usepackage{bitpolitostudiotemplates}
\usepackage{bitpolitostudiolibrary}
\usepackage[backend=biber]{biblatex}
\addbibresource{bibliography.bib}
\addbibresource{library/references.bib}
\begin{document}
\studiotemplate[
  palette={light},
  title={Privacy\\levels},
  item-one={
    text={Full-KYC: identity verification using submitted documents}
  },
  item-two={
    text={Soft-KYC: limited personal information is requested}
  },
  item-three={
    text={No-KYC: no personal information is requested}
  }
]{levels}
\studiotemplate[
  palette={blue},
  title={Privacy\\levels},
  item-one={
    text={Full-KYC: identity verification using submitted documents}
  },
  item-two={
    text={Soft-KYC: limited personal information is requested}
  },
  item-three={
    text={No-KYC: no personal information is requested}
  }
]{levels}
\studiotemplate[
  palette={light},
  title={Shared\\rules},
  item-one={
    text={Review the shared rules and compare the checks performed locally.}
  },
  item-two={
    text={Record the assumptions used by each participant in the example.}
  },
  item-three={
    text={Compare the available options.\\Explain the practical trade-offs.}
  }
]{levels}
\studiotemplate[
  palette={blue},
  title={Shared\\rules},
  item-one={
    text={Review the shared rules and compare the checks performed locally.}
  },
  item-two={
    text={Record the assumptions used by each participant in the example.}
  },
  item-three={
    text={Compare the available options.\\Explain the practical trade-offs.}
  }
]{levels}
\studiotemplate[
  palette={light},
  title={Icon\\choices},
  item-one={
    text={Check the documented rules.},
    icon={shield}
  },
  item-two={
    text={Study the example implementation.},
    icon={code}
  },
  item-three={
    text={Compare the available learning material.},
    icon={book}
  }
]{levels}
\studiotemplate[
  palette={blue},
  title={Icon\\choices},
  item-one={
    text={Check the documented rules.},
    icon={shield}
  },
  item-two={
    text={Study the example implementation.},
    icon={code}
  },
  item-three={
    text={Compare the available learning material.},
    icon={book}
  }
]{levels}
\studiotemplate[
  palette={light},
  title={Custom\\graphics},
  item-one={
    text={A tall graphic keeps its proportions.},
    graphic={\rule{12\studiounit}{220\studiounit}}
  },
  item-two={
    text={A wide graphic stays within the same bounds.},
    graphic={\rule{600\studiounit}{25\studiounit}}
  },
  item-three={
    text={An empty icon leaves the text in place.},
    icon={}
  }
]{levels}
\studiotemplate[
  palette={light},
  title={Your\\images},
  item-one={
    text={Your first image and description.},
    image={assets/generated/book.pdf}
  },
  item-two={
    text={Your second image and description.},
    image={assets/generated/code.pdf}
  },
  item-three={
    text={Your third image and description.},
    image={assets/generated/shield.pdf}
  }
]{levels}
\studiotemplate[
  palette={blue},
  title={Your\\images},
  item-one={
    text={Your first image and description.},
    image={images/studio/icon-palettes/book.pdf}
  },
  item-two={
    text={Your second image and description.},
    image={images/studio/icon-palettes/code.pdf}
  },
  item-three={
    text={Your third image and description.},
    image={images/studio/icon-palettes/shield.pdf}
  }
]{levels}
\studiotemplate[
  palette={light},
  title={},
  item-one={
    text={Keep this description beside its own icon.},
    icon={shield}
  },
  item-two={
    text={},
    icon={book}
  },
  item-three={
    text={A caption can remain without its icon.},
    icon={}
  }
]{levels}
\end{document}
